
\documentclass[11pt]{article}
\usepackage[utf8]{inputenc}
\usepackage[spanish]{babel}
\usepackage{amsmath}
\usepackage{geometry}
\geometry{margin=1in}

\title{Documentación: Ruido Blanco}
\date{}

\begin{document}
\maketitle

\section{Objetivo}
Analizar el concepto de ruido blanco dentro de las series de tiempo.

\section{Definición}
Un proceso de ruido blanco cumple:
\[
X_t \sim WN(0,\sigma^2)
\]

donde las observaciones son independientes, tienen media constante y varianza constante.

\section{Propiedades}
\begin{itemize}
\item Media igual a cero.
\item Varianza constante.
\item Ausencia de autocorrelación.
\item Comportamiento aleatorio.
\end{itemize}

\section{Importancia}
Los residuos de un buen modelo de series de tiempo deberían comportarse como ruido blanco.

\end{document}
